
\newpage
\section{System Overview}
\label{sec:baseline}

% Design principle for Chapter 3 > Describe the RAG system **as implemented**, prior to any optimization or evaluation.
% This chapter is strictly **descriptive**: no results, no comparison, no judgment.
% Reader takeaway after Chapter 3 > "I understand what the system is, what components and modules it contains, and how a query is executed in the baseline (always-on) configuration."

% ---------------------------------------------------------------------------
\subsection{The CITI KMS Context}
\label{subsec:kms-context}

The RAG system under study was deployed as a knowledge management system (KMS)
developed by the Center for IoT Innovation known as CITI, a leading research lab directed by Distinguished Professor 
Shuo-Yan Chou at the National Taiwan University of Science and Technology. 
While the system's primary purpose 
is to preserve the lab's institutional knowledge across student generations, 
it also serves as a first step toward what might be described as an organizational 
memory, with potential future use in internal agent-based 
workflows~\cite{openaidataagent2026}.

This chapter describes the system as physically deployed, focusing 
exclusively on its inference-time capabilities. Its special capabilities such as the podcast or teacher-student 
mode are not discussed.

% ---------------------------------------------------------------------------
\subsection{System Architecture}
\label{subsec:architecture}
The pipeline is deployed across four services: three primary processing 
nodes and one orchestrator.

\begin{figure}[h!]
    \centering
        \centering
        \includegraphics[width=0.85\textwidth]{nodes.png}
        \caption{CITI KMS infrastructure nodes}
        \label{fig:rag-vanilla}
\end{figure}

\begin{itemize}
    \item \textbf{Orchestrator} This service is the entry point for all user queries. It coordinates calls to the 
services, assembles intermediate outputs and returns the final answer. It performs no 
heavy computation. Its role is control flow and data assembly.
    \item \textbf{Embedding node} It contains the embedding model which encodes the text into dense vector representations. It processes user queries and 
when query expansion is active, the new generated query into the same vector space used for 
retrieval.
    \item \textbf{Database node} This node stores indexed document representations and executes similarity searches against 
incoming query vectors. It returns a ranked list of candidate passages to be used as context 
for generation.
    \item \textbf{LLM node} This node generates the final natural-language answer. It receives a prompt containing 
the user query together with relevant document passages retrieved from the database.
\end{itemize}



The Table below summarizes the hardware of each node. The orchestrator runs 
on a development machine without GPU.

\begin{table}[h!]
\centering
\begin{tabular}{llllll}
\textbf{Node} & \textbf{CPU} & \textbf{RAM} & \textbf{GPU}\\
\hline
LLM          & AMD Ryzen 9 9950X3D & 64\,GB & NVIDIA RTX 5090 (32\,GB) \\
Embedding   &  Intel i7-13700KF    & 32\,GB & NVIDIA RTX 4080 (16\,GB) \\
Database    &  Intel i7-12700      & 64\,GB & NVIDIA RTX 4080 (16\,GB) \\
Orchestrator &  Apple M4 Pro         & 24\,GB & ---                      \\

\end{tabular}
\caption{Node hardware specifications.}
\label{tab:deployment}
\end{table}

All nodes are connected via a local network. Energy measurements are collected on the 
three main nodes as described in Chapter~\ref{sec:method_protocol}.

% % Block C — Data flow at inference time
At inference time, the orchestrator drives a sequential flow: a user query arrives, 
the orchestrator requests a vector embedding, executes retrieval against the vector 
database, assembles the retrieved passages into a generation context and forwards it 
to the language model, which produces the response. Each stage has explicit inputs and 
outputs; the pipeline is modular by construction.

% \begin{figure}[h!]
%     \centering
%         \centering
%         \includegraphics[width=0.85\textwidth]{simple-rag.png}
%         \caption{Basic RAG pipeline with node attribution}
%         \label{fig:rag-simple-color}
% \end{figure}


% Block D — Architectural assumptions relevant to later chapters
For the purposes of this thesis, each request is treated as an independent unit of analysis: 
one query produces one execution trace. All experiments assume sequential execution with no 
concurrent requests.

% ---------------------------------------------------------------------------

\subsection{Optional Pipeline Modules}
\label{subsec:modules}

The CITI KMS extends the chore pipeline with three optional modules (introduced in Section~\ref{sec:architecture}):
two additional stages and one modification to an existing stage. Each module operates as
an independent toggle.

\medskip
\noindent\textit{Terminology note.}
Throughout this thesis, \emph{module} refers to a toggleable pipeline component; 
HyDE, reranking, hybrid search. \emph{Feature} refers to a predictive attribute
extracted from the query text (e.g.\ token count, embedding norm) used as input
to the router. This distinction avoids ambiguity in later chapters where both
concepts co-occur.


\paragraph{HyDE --- query expansion (before retrieval).}
When enabled, the orchestrator forwards the user query to the LLM node, which generates
a hypothetical answer. The embedding node then encodes this hypothetical document instead
of the original query. The resulting vector serves as an alternative query representation
for retrieval, bridging the vocabulary gap between short user queries and stored
documents. HyDE is the most expensive optional module: it adds one full LLM inference
call on the LLM node and one additional embedding call on the embedding node before
retrieval begins.

\begin{figure}[h!]
    \centering
    \includegraphics[width=.9\textwidth]{hyde-comparison.png}
    \caption{Query ingestion with and without HyDE}
    \label{fig:hyde-comparison}
\end{figure}

\paragraph{Hybrid retrieval --- search strategy (during retrieval).}
When enabled, the database node performs both dense (embedding-based) and sparse
(keyword-based) retrieval and fuses their ranked results. Dense search captures semantic
similarity; sparse search captures exact keyword matches. Combining both signals improves
coverage when one strategy alone would miss relevant passages. This computation runs
entirely on the database node. When disabled, only dense retrieval is used.

\begin{figure}[h!]
    \centering
    \includegraphics[width=.9\textwidth]{hybrid-dense.png}
    \caption{Dense versus hybrid retrieval}
    \label{fig:search-type}
\end{figure}


\paragraph{Reranking --- candidate refinement (after retrieval).}
When enabled, the embedding node runs a cross-encoder model over the retrieved candidate
documents. Unlike the embedding model, which encodes query and document independently, the
cross-encoder processes each query--document pair jointly and produces a relevance score.
The candidates are reordered by this score and only the top-$n$ passages are kept for
generation. Reranking adds one additional model call on the embedding node.

\begin{figure}[h!]
    \centering
    \includegraphics[width=0.9\textwidth]{reranking-comparison.png}
    \caption{Retrieved context with and without reranking}
    \label{fig:reranking}
\end{figure}

\medskip

Each module is an independent binary toggle. Together, the three toggles define
a configuration vector: HyDE (on/off), retrieval mode (dense/hybrid) and reranking
(on/off). This vector controls pipeline intensity per request. All three add computation and
service calls without changing the system's goal: producing an answer grounded in
retrieved context. During baseline experiments, the configuration is fixed for the
entire run and applied uniformly to all queries. Dynamic per-query adaptation is
introduced in Chapter~\ref{sec:controller}.


% ---------------------------------------------------------------------------
\subsection{Baseline Execution}
\label{subsec:baseline-flow}

In the baseline configuration, all optional modules are enabled. Every query follows the same 
fixed pipeline execution path with no conditional logic:

\begin{figure}[h!]
    \centering
        \centering
        \includegraphics[width=1.0\textwidth]{rag-options.png}
        \caption{RAG Full}
        \label{fig:flow-rag-hyde}
\end{figure}


\begin{enumerate}
    \item \textbf{Query received.} The orchestrator receives the user query.
    \item \textbf{HyDE generation.} The query is sent to the LLM node, which generates a hypothetical answer.
    \item \textbf{Embedding.} The embedding node encodes either the original query or the hypothetical document produced by HyDE.
    \item \textbf{Hybrid retrieval.} The database node performs dense and sparse search, fuses results and returns candidate documents.
    \item \textbf{Reranking.} The embedding node reranks the candidates and selects the top-$n$ passages.
    \item \textbf{Generation.} The LLM node generates the final answer conditioned on the reranked context.
    \item \textbf{Response returned.} The orchestrator returns the answer to the user.
\end{enumerate}

This execution is serial: each stage must complete before the next begins. All queries incur the full computational path regardless of their complexity or characteristics.

% ---------------------------------------------------------------------------
\subsection{Configuration Space}
\label{subsec:baseline-config}
The three optional modules described above: HyDE, reranking and retrieval
type. Each admits two states: on or off (dense or hybrid for retrieval type). 
Combined possibilities listed in Table~\ref{tab:config-space}.

% Combining
% all possibilities yields $2^3 = 8$ configurations, listed in
% Table~\ref{tab:config-space}.

\begin{table}[H]
\centering
\begin{tabular}{lccc}
\hline
\textbf{Config} & \textbf{HyDE} & \textbf{Reranking} & \textbf{Search type} \\
\hline
C0 & off & off & dense   \\
C1 & on  & off & dense   \\
C2 & off & on  & dense   \\
C3 & off & off & hybrid  \\
C4 & on  & on  & dense   \\
C5 & on  & off & hybrid  \\
C6 & off & on  & hybrid  \\
C7 & on  & on  & hybrid  \\
\hline
\end{tabular}
\caption{Configuration space ($2^3$ factorial design).}
\label{tab:config-space}
\end{table}

C0 serves as the minimal baseline: only the core pipeline stages execute. C7 is the
current CITI KMS deployment configuration and serves as the always-on reference against
which energy savings are measured. Single-module configurations (C1, C2, C3) enable
per-module attribution; multi-module configurations (C4--C6) test module interactions.
The following chapter describes the measurement framework used to quantify
the energy and performance characteristics across all eight configurations.
